%!TEX program = xelatex
% =====================================================================
%  COMPTE RENDU DE TRAVAUX PRATIQUES (TP)
%  Template LaTeX adapte du template de rapport de stage ENSIASD.
%  Compilateur : XeLaTeX. Compiler 2 a 3 fois (sommaire + renvois).
% =====================================================================
\documentclass[12pt,a4paper,oneside]{report}

% ---------------------------------------------------------------------
%  1. INFORMATIONS PERSONNALISABLES  (a modifier)
% ---------------------------------------------------------------------
\newcommand{\ecoleNom}{École Nationale Supérieure de l'Intelligence
  Artificielle et Sciences des Données}
\newcommand{\ecoleVille}{ENSIASD Taroudant}
\newcommand{\filiere}{Intelligence Artificielle et Sciences des Données}

\newcommand{\moduleNom}{Nom du module}          % ex: Programmation Python
\newcommand{\tpNumero}{TP N\textsuperscript{o}~1}
\newcommand{\tpTitre}{Titre du compte rendu de TP}

\newcommand{\etudiantNom}{Prénom Nom}
\newcommand{\etudiantGroupe}{Groupe ...}
\newcommand{\enseignantNom}{Pr. Nom de l'enseignant}
\newcommand{\dateRemise}{Le 1\textsuperscript{er} Janvier 2026}
\newcommand{\anneeUniv}{2025 / 2026}

% Logo de l'école (nom sans extension, présent dans images/)
\newcommand{\logoEcole}{images/logo_ensiasd}
% ---------------------------------------------------------------------

% ---------------------------------------------------------------------
%  2. PREAMBULE
% ---------------------------------------------------------------------
\usepackage[a4paper,top=2.5cm,bottom=2.5cm,left=2.5cm,right=2.5cm]{geometry}
\usepackage{amsmath,amssymb}
\usepackage{xcolor}

% --- Couleurs du document et de la coloration syntaxique ---
\definecolor{couvbleu}{RGB}{41,84,150}
\definecolor{codebg}{RGB}{246,248,251}
\definecolor{codekw}{RGB}{0,0,200}
\definecolor{codecomment}{RGB}{96,139,78}
\definecolor{codestring}{RGB}{163,21,21}
\definecolor{codenum}{RGB}{148,148,148}

\usepackage{graphicx}
\usepackage{tikz}
\usepackage{caption}
\usepackage{newfloat}
\usepackage{titlesec}
\usepackage{fancyhdr}
\usepackage{tabularx}
\usepackage{array}
\usepackage{booktabs}
\usepackage{ragged2e}
\usepackage{enumitem}
\usepackage{etoolbox}
\usepackage{url}

% --- Interligne 1,5 ---
\usepackage{setspace}
\onehalfspacing

% --- Evite les debordements de ligne (Overfull \hbox) ---
\setlength{\emergencystretch}{2em}

% --- Polices (TeX Gyre Termes ~ Times, TeX Gyre Cursor pour le code) ---
\usepackage{fontspec}
\setmainfont{TeX Gyre Termes}
\setmonofont{TeX Gyre Cursor}[Scale=0.9]

% --- Langue française (guillemets, typographie) ---
\usepackage{polyglossia}
\setdefaultlanguage{french}

% --- Coloration syntaxique + boîtes de code ---
\usepackage{listings}
\usepackage[most]{tcolorbox}

% =====================================================================
%  3. STYLE DE COLORATION SYNTAXIQUE (listings)
% =====================================================================
\lstdefinestyle{stylecode}{
  basicstyle=\ttfamily\small,
  commentstyle=\color{codecomment}\itshape,
  keywordstyle=\color{codekw}\bfseries,
  stringstyle=\color{codestring},
  numberstyle=\tiny\color{codenum},
  numbers=left,
  numbersep=10pt,
  breaklines=true,
  postbreak=\mbox{\textcolor{red}{$\hookrightarrow$}\space},
  showstringspaces=false,
  tabsize=2,
  columns=fullflexible,
  keepspaces=true,
  captionpos=b,
}
\lstset{style=stylecode}

% --- Style pour le code en ligne ---
\lstdefinestyle{styleinline}{basicstyle=\ttfamily\small,columns=fixed}
% Usage : \codeinline{print("bonjour")}
\newcommand{\codeinline}[1]{\colorbox{codebg}{\lstinline[style=styleinline]`#1`}}

% =====================================================================
%  4. BOITES DE CODE REUTILISABLES
% =====================================================================
% -- Boite de code avec titre (option : language=Python, numbers=left ...)
%    Usage : \begin{code}[language=Python]{bonjour.py} ... \end{code}
\newtcblisting{code}[2][]{%
  enhanced, breakable,
  listing only,
  listing engine=listings,
  colback=codebg,
  colframe=couvbleu,
  boxrule=0.5pt,
  arc=1.5mm, outer arc=1.5mm,
  left=3mm, right=3mm, top=2mm, bottom=2mm,
  toptitle=1.5mm, bottomtitle=1.5mm,
  fonttitle=\bfseries\small,
  coltitle=white,
  title={\ttfamily\detokenize{#2}},
  listing options={style=stylecode,#1}
}

% -- Boite de terminal / sortie console (fond sombre)
\lstdefinestyle{styleterm}{
  basicstyle=\ttfamily\small\color{white},
  numbers=none,
  breaklines=true,
  showstringspaces=false,
  tabsize=2,
  columns=fullflexible,
}
\newtcblisting{terminal}[1][]{%
  enhanced, breakable,
  listing only,
  listing engine=listings,
  colback=black!88,
  colframe=black,
  coltext=white!92,
  boxrule=0pt,
  arc=1.5mm, outer arc=1.5mm,
  left=3mm, right=3mm, top=2mm, bottom=2mm,
  listing options={style=styleterm,#1}
}

% -- Boite "remarque / consigne" (texte normal)
\newtcolorbox{remarque}[1][Remarque]{%
  enhanced, breakable,
  colback=blue!5,
  colframe=blue!55!black,
  boxrule=0.4pt,
  arc=1.5mm, outer arc=1.5mm,
  left=3mm, right=3mm, top=2mm, bottom=2mm,
  fonttitle=\bfseries,
  coltitle=white,
  title={#1}
}

% =====================================================================
%  5. TITRES DES PARTIES ET SECTIONS
% =====================================================================
% Parties (chapitres) : "Partie I : TITRE", 16pt gras, aligné à gauche
\titleformat{\chapter}[block]
  {\normalfont\bfseries}
  {\fontsize{16}{20}\selectfont\MakeUppercase{Partie \thechapter}~:}
  {0.4em}
  {\raggedright\normalfont\fontsize{16}{20}\selectfont\bfseries\MakeUppercase}
\titleformat{name=\chapter,numberless}[block]
  {\normalfont\bfseries}
  {}
  {0pt}
  {\raggedright\normalfont\fontsize{16}{20}\selectfont\bfseries\MakeUppercase}
\titlespacing*{\chapter}{0pt}{0pt}{54pt}
\titlespacing*{name=\chapter,numberless}{0pt}{0pt}{54pt}

\renewcommand{\thechapter}{\Roman{chapter}}

% Sections : "I.1 Titre", 14pt gras
\titleformat{\section}
  {\normalfont\fontsize{14}{17}\selectfont\bfseries}
  {\thesection}
  {0.6em}
  {}
\titlespacing*{\section}{0pt}{12pt}{12pt}

% Sous-sections : "I.1.1 Titre", 12pt gras
\titleformat{\subsection}
  {\normalfont\fontsize{12}{15}\selectfont\bfseries}
  {\thesubsection}
  {0.6em}
  {}
\titlespacing*{\subsection}{0pt}{6pt}{6pt}

% Subdivisions : "I.1.1.1 Titre", 12pt gras
\titleformat{\subsubsection}
  {\normalfont\fontsize{12}{15}\selectfont\bfseries}
  {\thesubsubsection}
  {0.6em}
  {}
\titlespacing*{\subsubsection}{0pt}{6pt}{6pt}

\renewcommand{\thesection}{\thechapter.\arabic{section}}
\renewcommand{\thesubsection}{\thesection.\arabic{subsection}}
\renewcommand{\thesubsubsection}{\thesubsection.\arabic{subsubsection}}

\setcounter{secnumdepth}{3}

% =====================================================================
%  6. FIGURES / TABLEAUX
% =====================================================================
\renewcommand{\thefigure}{\thechapter.\arabic{figure}}
\renewcommand{\thetable}{\thechapter.\arabic{table}}

\captionsetup[figure]{
  font={small,it}, justification=centering,
  labelfont=bf, labelsep=space, skip=6pt, position=bottom
}
\captionsetup[table]{
  font=small, justification=centering,
  labelfont=bf, labelsep=space, skip=6pt, position=top
}
\newcommand{\source}[1]{\par\footnotesize\itshape Source~: #1}

% =====================================================================
%  7. EN-TETES ET PIEDS DE PAGE
% =====================================================================
\pagestyle{fancy}
\fancyhf{}
\fancyfoot[R]{\thepage}
\renewcommand{\headrulewidth}{0pt}
\fancypagestyle{plain}{%
  \fancyhf{}
  \fancyfoot[R]{\thepage}
  \renewcommand{\headrulewidth}{0pt}%
}

% =====================================================================
%  8. LIENS HYPERTEXTE (chargé en dernier)
% =====================================================================
\usepackage[colorlinks=true, linkcolor=couvbleu, citecolor=couvbleu,
  urlcolor=couvbleu]{hyperref}

% --- Bande bleue pleine largeur (page de garde) ---
\newcommand{\bandebleue}[1]{%
  \par\noindent
  \hspace*{-\dimexpr\oddsidemargin+1in\relax}%
  \textcolor{couvbleu}{\rule{\paperwidth}{#1}}%
  \par
}

% =====================================================================
%  9. DEBUT DU DOCUMENT
% =====================================================================
\begin{document}

% --- Page de garde (non numérotée) ---
\pagestyle{empty}
% =====================================================================
%  PAGE DE GARDE DU COMPTE RENDU DE TP  (front/cover.tex)
% =====================================================================
\thispagestyle{empty}

% --- Logo seul, centré ---
\vspace*{-1cm}
\begin{center}
  \IfFileExists{\logoEcole.png}{%
    \includegraphics[width=0.55\textwidth]{\logoEcole.png}}{%
    \IfFileExists{\logoEcole.jpg}{%
      \includegraphics[width=0.55\textwidth]{\logoEcole.jpg}}{%
      \textcolor{gray}{[logo]}}}
\end{center}

% --- Nom de l'école ---
\vspace{0.2cm}
\begin{center}
  {\small\ecoleNom}
\end{center}

% --- Bande bleue ---
\vspace{0.3cm}
\bandebleue{0.32cm}
\vspace{0.7cm}

\begin{spacing}{1.25}
\begin{center}
  {\fontsize{20}{24}\selectfont\bfseries
    \MakeUppercase{Compte Rendu de Travaux Pratiques}}

  \vspace{0.5cm}
  {\bfseries Module~: «~\moduleNom~»}

  \vspace{0.7cm}
  \textcolor{couvbleu}{\rule{\linewidth}{1.2mm}}

  \vspace{0.4cm}
  {\Large\bfseries \tpNumero}
  \vspace{0.2cm}

  {\fontsize{18}{22}\selectfont\bfseries \tpTitre}

  \vspace{0.4cm}
  \textcolor{couvbleu}{\rule{\linewidth}{1.2mm}}
\end{center}
\end{spacing}

% --- Réalisé par / Encadré par ---
\vspace{1cm}
\begin{spacing}{1.2}
\noindent
\begin{tabularx}{\linewidth}{@{}>{\bfseries}p{4.5cm}X@{}}
  Filière~:        & \filiere \\[6pt]
  Réalisé par~:    & \etudiantNom \\[6pt]
  Groupe~:         & \etudiantGroupe \\[6pt]
  Encadré par~:    & \enseignantNom \\[6pt]
  Date de remise~: & \dateRemise \\
\end{tabularx}
\end{spacing}

% --- Année universitaire (bas de page) ---
\vfill
\begin{center}
  \textbf{Année Universitaire~: \anneeUniv}
\end{center}
\vspace{0.35cm}
\bandebleue{0.32cm}
\vspace*{0.2cm}

\clearpage


% --- Sommaire ---
\clearpage
\pagestyle{plain}
\pagenumbering{Roman}
\phantomsection
\addcontentsline{toc}{chapter}{Sommaire}
\tableofcontents
\clearpage

% --- Listes optionnelles : décommenter si besoin ---
% \phantomsection
% \addcontentsline{toc}{chapter}{Liste des figures}
% \listoffigures
% \clearpage

% =====================================================================
%  CORPS DU DOCUMENT
% =====================================================================
\pagenumbering{arabic}
\setcounter{page}{1}
\pagestyle{fancy}
\fancyhf{}
\fancyfoot[R]{\thepage}
\renewcommand{\headrulewidth}{0pt}
\fancypagestyle{plain}{%
  \fancyhf{}
  \fancyfoot[R]{\thepage}
  \renewcommand{\headrulewidth}{0pt}%
}

% =====================================================================
%  INTRODUCTION, OBJECTIFS ET ENVIRONNEMENT DE TRAVAIL (non numéroté)
% =====================================================================
\chapter*{Introduction et objectifs}
\addcontentsline{toc}{chapter}{Introduction et objectifs}

Ce document constitue le compte rendu des travaux pratiques réalisés dans le
cadre du module «~\moduleNom~». Il présente la démarche suivie, le code
développé ainsi que les résultats obtenus et leur interprétation.

\paragraph{Objectifs du TP.}
À l'issue de cette séance, l'étudiant doit être capable de~:
\begin{itemize}[leftmargin=1.2em,itemsep=1pt]
  \item citer ici le premier objectif pédagogique du TP ;
  \item manipuler les outils et concepts abordés (langage, bibliothèque,
    algorithme\ldots) ;
  \item analyser et interpréter les résultats produits par le programme.
\end{itemize}

\paragraph{Environnement de travail.}
Indiquez ici le matériel et les logiciels utilisés~:
\begin{itemize}[leftmargin=1.2em,itemsep=1pt]
  \item \textbf{Système d'exploitation}~: Windows 11 / Linux / macOS ;
  \item \textbf{Langage / version}~: \codeinline{Python 3.12}, \codeinline{gcc 13},
    \codeinline{Java 21}\ldots;
  \item \textbf{Éditeur / IDE}~: VS Code, PyCharm, IntelliJ, Vim\ldots;
  \item \textbf{Outils complémentaires}~: Git, terminal, Jupyter Notebook.
\end{itemize}

\paragraph{Organisation du rapport.}
La \textbf{Première partie} présente le premier exercice traité. La
\textbf{Deuxième partie} détaille un second exercice. Une conclusion générale
clôt le document.

\begin{remarque}[Comment utiliser ce template]
Remplacez le texte de chaque partie par votre propre contenu. Les blocs de
code s'écrivent avec l'environnement \texttt{code} (options
\texttt{language=\ldots} et titre entre accolades), et les sorties de terminal
avec l'environnement \texttt{terminal}. Voir le fichier \texttt{README.md}
pour la liste des langages supportés et des exemples complets.
\end{remarque}

% =====================================================================
%  PARTIE I : exemple d'exercice avec du code Python
% =====================================================================
\chapter{Prise en main de Python et structures de contrôle}

\section{Énoncé de l'exercice}

\begin{remarque}[Consigne]
Écrire un programme Python qui calcule la moyenne d'une liste de notes, puis
affiche la mention correspondante en utilisant la structure conditionnelle
\codeinline{if / elif / else}.
\end{remarque}

Le programme doit demander à l'utilisateur de saisir les notes, calculer la
moyenne et afficher un message indiquant si l'étudiant est admis.

\section{Démarche et implémentation}

Nous définissons d'abord une fonction \codeinline{moyenne} qui prend une liste
en paramètre et renvoie sa moyenne. La fonction \codeinline{mention} associe
ensuite une mention à la moyenne obtenue.

\begin{code}[language=Python]{notes.py}
# Calcul d'une moyenne et affichage de la mention associee

def moyenne(notes):
    """Retourne la moyenne d'une liste de notes."""
    if not notes:
        return 0.0
    return sum(notes) / len(notes)


def mention(moy):
    """Retourne la mention correspondant a une moyenne."""
    if moy >= 16:
        return "Tres bien"
    elif moy >= 14:
        return "Bien"
    elif moy >= 12:
        return "Assez bien"
    elif moy >= 10:
        return "Admis"
    else:
        return "Ajourne"


if __name__ == "__main__":
    notes = [15.5, 12.0, 17.0, 9.5]
    moy = moyenne(notes)
    print(f"Moyenne : {moy:.2f}")
    print(f"Mention : {mention(moy)}")
\end{code}

\section{Résultats et interprétation}

L'exécution du programme dans le terminal produit la sortie suivante.

\begin{terminal}
$ python notes.py
Moyenne : 13.50
Mention : Assez bien
\end{terminal}

La moyenne calculée vaut \textbf{13{,}50}, ce qui correspond à la mention
«~Assez bien~». Le test de la fonction \codeinline{mention} avec différentes
valeurs permet de vérifier que chaque branche conditionnelle renvoie la
mention attendue (voir Tableau~\ref{tab:mentions}).

\begin{table}[h]
  \centering
  \begin{tabularx}{0.75\linewidth}{@{}Xc@{}}
    \toprule
    \textbf{Intervalle de la moyenne} & \textbf{Mention} \\
    \midrule
    $[16~;~20]$   & Très bien  \\
    $[14~;~16[$   & Bien       \\
    $[12~;~14[$   & Assez bien \\
    $[10~;~12[$   & Admis      \\
    $[0~;~10[$    & Ajourné    \\
    \bottomrule
  \end{tabularx}
  \caption{Correspondance entre moyenne et mention.}
  \label{tab:mentions}
\end{table}

\begin{remarque}[Interprétation]
Le programme respecte la consigne. On remarque que l'absence de notes est
gérée par le cas de base \codeinline{not notes} afin d'éviter une division par
zéro. Toutes les valeurs limites ont été testées.
\end{remarque}

% =====================================================================
%  PARTIE II : exemple multi-langages (C, SQL) + figure
% =====================================================================
\chapter{Tri d'un tableau et interrogation d'une base de données}

\section{Énoncé de l'exercice}

\begin{remarque}[Consigne]
Écrire un programme C qui trie un tableau d'entiers par la méthode du tri à
bulles, puis une requête SQL qui sélectionne les employés dont le salaire
dépasse un seuil donné.
\end{remarque}

\section{Implémentation en langage C}

Le tri à bulles parcourt le tableau et échange deux éléments consécutifs
lorsqu'ils sont mal ordonnés. Nous utilisons la fonction \codeinline{swap}
pour l'échange.

\begin{code}[language=C, numbers=left]{tri.c}
#include <stdio.h>
#include <stddef.h>

static void swap(int *a, int *b)
{
    int tmp = *a;
    *a = *b;
    *b = tmp;
}

void tri_bulles(int *tab, size_t n)
{
    for (size_t i = 0; i < n - 1; ++i) {
        for (size_t j = 0; j < n - 1 - i; ++j) {
            if (tab[j] > tab[j + 1]) {
                swap(&tab[j], &tab[j + 1]);
            }
        }
    }
}

int main(void)
{
    int tab[] = {5, 2, 9, 1, 7};
    size_t n = sizeof(tab) / sizeof(tab[0]);

    tri_bulles(tab, n);

    for (size_t i = 0; i < n; ++i) {
        printf("%d ", tab[i]);
    }
    printf("\n");
    return 0;
}
\end{code}

Compilation et exécution~:

\begin{terminal}
$ gcc -Wall -o tri tri.c
$ ./tri
1 2 5 7 9
\end{terminal}

\section{Requête SQL associée}

\begin{code}[language=SQL]{employes.sql}
SELECT nom, prenom, salaire
FROM employes
WHERE salaire > 8000
ORDER BY salaire DESC;
\end{code}

Le résultat de la requête peut être présenté sous forme de capture d'écran
(Figure~\ref{fig:capture}).

\begin{figure}[h]
  \centering
  \includegraphics[width=0.7\linewidth]{images/capture_exemple.png}
  \caption{Exemple de résultat obtenu après exécution.}
  \label{fig:capture}
\end{figure}

\section{Analyse et complexité}

Le tri à bulles a une complexité temporelle en $O(n^2)$ dans le pire cas, ce
qui le rend peu efficace pour de grands tableaux. La requête SQL, quant à
elle, s'appuie sur l'opérateur de comparaison \codeinline{salaire > 8000} et
sur la clause \codeinline{ORDER BY} pour trier le résultat.

% =====================================================================
%  CONCLUSION (non numéroté)
% =====================================================================
\chapter*{Conclusion}
\addcontentsline{toc}{chapter}{Conclusion}

Ce TP a permis de mettre en pratique les notions essentielles du module
«~\moduleNom~»~: manipulation des structures de données, écriture et
exécution de programmes, puis analyse des résultats. Les objectifs fixés ont
été atteints et le code produit fonctionne correctement sur les jeux de
tests considérés.

Comme piste d'amélioration, on pourra~:
\begin{itemize}[leftmargin=1.2em,itemsep=1pt]
  \item optimiser l'algorithme de tri (tri rapide, tri fusion) ;
  \item ajouter une validation des entrées utilisateur ;
  \item écrire des tests unitaires pour automatiser la vérification.
\end{itemize}


\end{document}
